% Chapter 12: Generalization Bounds
%   uniform_convergence: Profile B (load-bearing theorem: symmetrization
%     + concentration proof IS the content)
%   Five bound traditions: Profile G (framework bridge: parallax)
%   rademacher_complexity: Profile A-to-B bridge (definition matters, but
%     comparison to VC is the content)
%   pac_bayes_bound: Profile E (obstruction: looks Bayesian but prior
%     has a fundamentally different role)
%
% Centerpiece: the symmetrization argument (Sec 2) and the five-framework
% parallax (Sec 8).

\chapter{Generalization Bounds}
\label{ch:generalization}

\Cref{ch:pac} established \emph{that} finite VC dimension characterizes PAC
learnability.  \Cref{ch:compression} refined the quantitative dependence
through compression.  This chapter asks a different question: given a
specific learning algorithm~$A$ and a specific training sample~$S$, what
can we say about the generalization error of the hypothesis~$A(S)$?

The answer is not unique.  Five distinct research traditions have developed
bounds on the same quantity (the gap between true and empirical
risk), each using different information about the learner:

\begin{enumerate}[leftmargin=*, itemsep=3pt]
\item \textbf{VC/uniform convergence} measures the combinatorial richness
  of the hypothesis class through the growth function.
\item \textbf{Rademacher complexity} refines this to a data-dependent
  measure of how well the class correlates with random noise.
\item \textbf{PAC-Bayes bounds} measure the ``distance'' of a learned
  posterior from a data-independent prior.
\item \textbf{Algorithmic stability} measures how much the output changes
  when a single training example is perturbed.
\item \textbf{Margin theory} measures the ``confidence'' of predictions
  and connects to scale-sensitive dimensions.
\end{enumerate}

\noindent
We also discuss information-theoretic bounds, which measure the mutual
information between the training sample and the learned hypothesis.

Each framework produces a bound on the same quantity but requires
different structural assumptions and succeeds in different regimes.  The
centerpiece of this chapter is twofold: (i)~the \emph{symmetrization
argument} that powers uniform convergence (\Cref{sec:uniform-convergence}),
proved in full as the engine behind classical generalization theory; and
(ii)~the \emph{comparison} among the five frameworks
(\Cref{sec:comparison}), which reveals what each sees that the others miss.

\medskip

Throughout this chapter, $\mathcal{H}$ denotes a hypothesis class,
$D$~a distribution on $X \times Y$,
$S = \{(x_i, y_i)\}_{i=1}^m$ an i.i.d.\ sample from~$D$,
$A$~a learning algorithm, and
$\ell \colon \mathcal{H} \times (X \times Y) \to [0,1]$ a bounded loss
function.


% ================================================================
% SECTION 1: GENERALIZATION ERROR
% Profile A.  Definition.
% ================================================================

\section{Generalization Error}
\label{sec:gen-error}

\begin{defn}[Generalization Error]
\label{def:gen-error}
The \emph{generalization error} (or \emph{generalization gap}) of a
hypothesis $h$ with respect to distribution~$D$ and sample~$S$ is
\[
  \mathrm{gen}(h, S) \;=\; \risk_D(h) - \emprisk_S(h)
    \;=\; \E_{(x,y) \sim D}[\ell(h, (x,y))]
    - \frac{1}{m}\sum_{i=1}^m \ell(h, (x_i, y_i)).
\]
A \emph{generalization bound} is any high-probability or in-expectation
upper bound on $|\mathrm{gen}(h, S)|$ or on $\mathrm{gen}(A(S), S)$.
\end{defn}

The generalization error is a random variable (through~$S$), and when
$h = A(S)$ depends on~$S$, controlling it requires accounting for the
dependence between the hypothesis and the data that produced it.  This is
the core technical challenge that each framework addresses differently.


% ================================================================
% SECTION 2: UNIFORM CONVERGENCE
% Profile B (load-bearing theorem).  FULL proof.
% ================================================================

\section{Uniform Convergence: The Engine}
\label{sec:uniform-convergence}

The classical approach to generalization does not look at the algorithm at
all.  Instead, it asks: does the empirical risk converge to the true risk
\emph{uniformly over all hypotheses}?  If so, any hypothesis with small
empirical risk automatically has small true risk, including the one the
algorithm selects.

\begin{defn}[Uniform Convergence]
\label{def:uniform-convergence}
A hypothesis class~$\mathcal{H}$ has the \emph{uniform convergence
property} with respect to loss~$\ell$ if, for every $\varepsilon > 0$,
\[
  \sup_{h \in \mathcal{H}} |\risk_D(h) - \emprisk_S(h)|
  \xrightarrow{\;m \to \infty\;} 0
  \quad \text{in probability, uniformly over all distributions } D.
\]
\end{defn}

The content of uniform convergence is the proof technique: the
\emph{symmetrization argument} introduced by Vapnik and Chervonenkis in
1971.  The argument proceeds in three stages, each with a distinct role.

\begin{thm}[VC Uniform Convergence Bound]
\label{thm:vc-uniform}
Let $\mathcal{H} \subseteq \{0,1\}^X$ with $\VC(\mathcal{H}) = d < \infty$.
Then for any distribution~$D$ on $X \times \{0,1\}$, any $\varepsilon > 0$,
and any $\delta > 0$, with probability at least $1 - \delta$ over
$S \sim D^m$:
\[
  \sup_{h \in \mathcal{H}} |\risk_D(h) - \emprisk_S(h)|
  \;\leq\; \sqrt{\frac{8\,d \log(2em/d) + 8\log(4/\delta)}{m}}.
\]
\end{thm}

\begin{proof}
The proof has three stages: symmetrization, reduction to a finite
problem, and concentration.

\textbf{Stage 1: Symmetrization (the ghost sample).}
Let $S = (z_1, \ldots, z_m)$ be the training sample and let
$S' = (z_1', \ldots, z_m')$ be an independent ``ghost'' sample drawn
from the same distribution.  Define
$\Phi(S) = \sup_{h \in \mathcal{H}} (\risk_D(h) - \emprisk_S(h))$.

For any fixed~$h$, the true risk equals
$\risk_D(h) = \E_{S'}[\emprisk_{S'}(h)]$, so
\[
  \Phi(S) = \sup_{h \in \mathcal{H}} \E_{S'}\!\left[
    \emprisk_{S'}(h) - \emprisk_S(h)\right].
\]
Moving the supremum inside the expectation (Jensen's inequality applied
to the convex functional $\sup$) gives:
\begin{equation}
\label{eq:symm-step}
  \E_S[\Phi(S)]
  \;\leq\; \E_{S,S'}\!\left[\sup_{h \in \mathcal{H}}
    \bigl(\emprisk_{S'}(h) - \emprisk_S(h)\bigr)\right].
\end{equation}

\emph{Why this helps.}  The left side involves the population
quantity~$\risk_D(h)$, which we cannot compute.  The right side involves
only \emph{two finite samples}; no population quantities remain.  This is
the symmetrization trick: the ghost sample replaces the unknown
distribution.

\textbf{Stage 1b: Rademacher sign-flip.}
The difference $\emprisk_{S'}(h) - \emprisk_S(h) =
\frac{1}{m}\sum_{i=1}^m (\ell(h,z_i') - \ell(h,z_i))$ is a sum of
symmetric random variables (since $z_i$ and $z_i'$ are identically
distributed, swapping them does not change the joint distribution).
Introducing i.i.d.\ Rademacher variables $\sigma_1,\ldots,\sigma_m$
(uniform on $\{-1,+1\}$):
\begin{align}
  \E_{S,S'}\!\left[\sup_{h} \frac{1}{m}\sum_{i=1}^m
    \bigl(\ell(h,z_i') - \ell(h,z_i)\bigr)\right]
  &= \E_{S,S',\sigma}\!\left[\sup_{h} \frac{1}{m}\sum_{i=1}^m
    \sigma_i\bigl(\ell(h,z_i') - \ell(h,z_i)\bigr)\right]
    \nonumber \\
  &\leq 2\,\E_{S,\sigma}\!\left[\sup_{h} \frac{1}{m}\sum_{i=1}^m
    \sigma_i \ell(h,z_i)\right],
    \label{eq:rademacher-step}
\end{align}
where the last step uses the triangle inequality for suprema and the
identical distribution of $S$ and~$S'$.

\textbf{Stage 2: Reduction to a finite class.}
The supremum in \eqref{eq:rademacher-step} ranges over all
$h \in \mathcal{H}$, but on the fixed sample~$S$, only finitely many
distinct behavior patterns occur.  Specifically, the loss vectors
$(\ell(h,z_1), \ldots, \ell(h,z_m))$ take at most $\growth_\mathcal{H}(m)$
distinct values, where $\growth_\mathcal{H}$ is the growth function.

For binary classification with 0--1 loss, each behavior pattern is a
binary vector in $\{0,1\}^m$, and the Sauer--Shelah lemma
(\Cref{ch:dimensions}) gives
\begin{equation}
\label{eq:sauer-shelah-app}
  \growth_\mathcal{H}(m) \;\leq\; \left(\frac{em}{d}\right)^d.
\end{equation}

Applying Massart's finite lemma to the at most $\growth_\mathcal{H}(m)$
distinct vectors, each with Euclidean norm at most
$\sqrt{m}$,\formal{FLT_Proofs.Complexity.Rademacher.finite_massart_lemma}
\begin{equation}
\label{eq:massart}
  \E_\sigma\!\left[\sup_{h} \frac{1}{m}\sum_{i=1}^m \sigma_i \ell(h,z_i)
  \right]
  \;\leq\; \sqrt{\frac{2\log \growth_\mathcal{H}(m)}{m}}
  \;\leq\; \sqrt{\frac{2d\log(em/d)}{m}}.
\end{equation}

\textbf{Stage 3: Concentration.}
The bound so far controls $\E[\Phi(S)]$.  To obtain a high-probability
bound, observe that $\Phi(S) = \sup_h (\risk_D(h) - \emprisk_S(h))$
satisfies the bounded-differences condition: replacing any single~$z_i$
changes $\Phi(S)$ by at most $1/m$ (since each loss value lies in
$[0,1]$).  By McDiarmid's inequality,
\[
  \Pr\!\left[\Phi(S) > \E[\Phi(S)] + t\right]
  \leq \exp(-2mt^2).
\]
Setting $t = \sqrt{\log(2/\delta)/(2m)}$ and combining with
\eqref{eq:rademacher-step}--\eqref{eq:massart}, the one-sided bound
follows.  Applying the same argument to
$\sup_h (\emprisk_S(h) - \risk_D(h))$ and taking a union bound yields
the two-sided result.
\end{proof}

\begin{historical}
The symmetrization argument was introduced by Vapnik and Chervonenkis
(1971) in the original paper that defined what is now called the VC
dimension.  The three-stage structure (ghost sample, sign-flip,
finite reduction) was not initially presented in this clean form.  The
modern presentation, which passes through Rademacher complexity as an
intermediate quantity, crystallized in the work of Koltchinskii (2001)
and Bartlett--Mendelson (2002).  The Rademacher step (Stage~1b) was
implicit in the original argument but was not isolated as a self-standing
concept until much later.
\end{historical}

\begin{traversal}
The uniform convergence theorem ties three threads together: the growth
function characterizes uniform convergence, the Sauer--Shelah lemma is the
step that turns this into the VC generalization bound, and the VC dimension
upper-bounds the Rademacher complexity.  The proof above traces exactly this
path: the growth function enters at Stage~2, Sauer--Shelah provides the
polynomial bound, and the Rademacher intermediate appears at Stage~1b.
\end{traversal}


\begin{example}[Halfspaces in $\R^d$]
\label{ex:halfspace-uc}
Let $\mathcal{H} = \{x \mapsto \ind[\langle w, x \rangle \geq b]
: w \in \R^d, b \in \R\}$ be the class of halfspaces.
Then $\VC(\mathcal{H}) = d+1$ (\Cref{ch:dimensions}), and the uniform
convergence bound gives: with probability $\geq 1 - \delta$,
\[
  \sup_{h \in \mathcal{H}} |\risk_D(h) - \emprisk_S(h)|
  \leq \sqrt{\frac{8(d+1)\log(2em/(d+1)) + 8\log(4/\delta)}{m}}.
\]
For $d = 10$ and $m = 10{,}000$, this gives a bound of approximately
$0.12$, a nontrivial and meaningful result.  For a two-layer neural network
with $p = 10^6$ parameters (where naive VC bounds give
$\VC \approx O(p \log p)$), the same formula gives a bound exceeding~$1$:
the uniform convergence bound is \emph{vacuous}.
\end{example}


% ================================================================
% SECTION 3: RADEMACHER COMPLEXITY
% Profile A-to-B bridge.  Definition + comparison to VC.
% ================================================================

\section{Rademacher Complexity}
\label{sec:rademacher}

The Rademacher intermediate that appeared in Stage~1b of the
symmetrization proof is itself a fundamental complexity measure.  It
refines VC dimension in two ways: it is \emph{data-dependent} (computed on
the actual sample, not worst-case) and it applies to \emph{real-valued}
function classes (not just binary hypotheses).

\begin{defn}[Empirical Rademacher Complexity]
\label{def:rademacher}
Let $\mathcal{F}$ be a class of functions $f \colon Z \to \R$ and let
$S = (z_1, \ldots, z_m)$ be a fixed sample.  The \emph{empirical
Rademacher complexity} of $\mathcal{F}$ on~$S$ is
\[
  \Rad_S(\mathcal{F}) = \E_{\sigma}\!\left[\sup_{f \in \mathcal{F}}
    \frac{1}{m}\sum_{i=1}^m \sigma_i f(z_i)\right],
\]
where $\sigma_1, \ldots, \sigma_m$ are independent Rademacher random
variables (uniform on $\{-1, +1\}$).  The \emph{Rademacher complexity} of
$\mathcal{F}$ at sample size~$m$ is
$\mathcal{R}_m(\mathcal{F}) = \E_S[\Rad_S(\mathcal{F})]$.\formal{FLT_Proofs.Complexity.Rademacher.EmpiricalRademacherComplexity}
\end{defn}

The Rademacher complexity measures how well $\mathcal{F}$ can
\emph{correlate with pure noise}.  A class that achieves high correlation
with random $\pm 1$ labels is rich enough to fit arbitrary patterns,
including the noise in the training data.

\begin{thm}[Rademacher Generalization Bound
{\cite{BartlettMendelson2002}}]
\label{thm:rademacher-bound}
Let $\mathcal{F}$ be a class of functions $f \colon Z \to [0,1]$, let
$D$ be a distribution on~$Z$, and let $S = (z_1,\ldots,z_m) \sim D^m$.
Then for any $\delta > 0$, with probability at least $1-\delta$ over~$S$,
for all $f \in \mathcal{F}$ simultaneously:
\[
  \E_D[f] \leq \frac{1}{m}\sum_{i=1}^m f(z_i)
    + 2\,\Rad_S(\mathcal{F})
    + 3\sqrt{\frac{\log(2/\delta)}{2m}}.
\]
\end{thm}

\begin{proof}
The proof follows the same three-stage architecture as the uniform
convergence theorem: symmetrization, Rademacher sign-flip, and
concentration.

Define $\Phi(S) = \sup_{f \in \mathcal{F}} (\E_D[f] -
\frac{1}{m}\sum_i f(z_i))$.

\textbf{Step 1: Concentration of $\Phi$.}
Replacing any single $z_i$ in~$S$ changes $\Phi(S)$ by at most $1/m$
(since $f$ is bounded in $[0,1]$).  By McDiarmid's inequality:
\begin{equation}
\label{eq:rad-conc}
  \Pr[\Phi(S) > \E[\Phi(S)] + t] \leq \exp(-2mt^2).
\end{equation}
Setting $t = \sqrt{\log(2/\delta)/(2m)}$ gives
$\Phi(S) \leq \E[\Phi(S)] + \sqrt{\log(2/\delta)/(2m)}$
with probability $\geq 1 - \delta/2$.

\textbf{Step 2: Symmetrization.}
Let $S' = (z_1', \ldots, z_m')$ be an independent ghost sample.  Then
\begin{align*}
  \E_S[\Phi(S)]
  &= \E_S\!\left[\sup_{f \in \mathcal{F}} \left(\E_{S'}
    \!\left[\frac{1}{m}\sum_{i=1}^m f(z_i')\right]
    - \frac{1}{m}\sum_{i=1}^m f(z_i)\right)\right] \\
  &\leq \E_{S, S'}\!\left[\sup_{f \in \mathcal{F}}
    \frac{1}{m}\sum_{i=1}^m \bigl(f(z_i') - f(z_i)\bigr)\right].
\end{align*}
Since $z_i$ and $z_i'$ are i.i.d., the differences are symmetric
random variables, so introducing Rademacher signs does not change the
distribution:
\begin{align*}
  \E_{S, S'}\!\left[\sup_{f}
    \frac{1}{m}\sum_i \bigl(f(z_i') - f(z_i)\bigr)\right]
  &= \E_{S, S', \sigma}\!\left[\sup_{f}
    \frac{1}{m}\sum_i \sigma_i\bigl(f(z_i') - f(z_i)\bigr)\right] \\
  &\leq 2\,\E_{S, \sigma}\!\left[\sup_{f}
    \frac{1}{m}\sum_i \sigma_i f(z_i)\right]
  = 2\,\mathcal{R}_m(\mathcal{F}).
\end{align*}

\textbf{Step 3: Data-dependent form.}
Apply McDiarmid again: $\Rad_S(\mathcal{F})$ has bounded differences
$1/m$ (changing one sample point changes the Rademacher average by at
most $1/m$), so
$\mathcal{R}_m(\mathcal{F}) \leq \Rad_S(\mathcal{F})
  + \sqrt{\log(2/\delta)/(2m)}$
with probability $\geq 1 - \delta/2$.

A union bound over the two $\delta/2$ events yields the stated bound.
\end{proof}


\subsection{Rademacher Complexity versus VC Dimension}
\label{sec:rad-vs-vc}

The following proposition makes the relationship precise.

\begin{prop}[VC dimension bounds Rademacher complexity]
\label{prop:rad-vc}
For any binary hypothesis class $\mathcal{H} \subseteq \{0,1\}^X$ with
$\VC(\mathcal{H}) = d$:
\[
  \mathcal{R}_m(\mathcal{H})
  \;\leq\; \sqrt{\frac{2d \log(em/d)}{m}}.
\]
This quantitative bound is machine-checked.\formal{FLT_Proofs.Complexity.Rademacher.vcdim_bounds_rademacher_quantitative}
\end{prop}

\begin{proof}
On any fixed sample of size~$m$, the restriction $\mathcal{H}|_S$ has at
most $\growth_\mathcal{H}(m) \leq (em/d)^d$ elements (Sauer--Shelah).
Each element is a binary vector of Euclidean norm at most $\sqrt{m}$.
By Massart's finite lemma, the Rademacher complexity of a finite set~$V$
of vectors in $\R^m$ satisfies
$\E_\sigma[\sup_{v \in V} \frac{1}{m}\langle \sigma, v\rangle]
\leq \frac{\max_v \|v\|_2}{m}\sqrt{2\log|V|}$.
Substituting $|V| \leq (em/d)^d$ and $\|v\|_2 \leq \sqrt{m}$ gives
the result.
\end{proof}

\begin{remark}[When Rademacher is tighter than VC]
\label{rmk:rad-tighter}
The VC bound is worst-case: it uses the maximum of
$|\mathcal{H}|_S|$ over all samples of size~$m$.  The empirical
Rademacher complexity $\Rad_S(\mathcal{H})$ uses the \emph{actual}
restriction $\mathcal{H}|_S$, which may be much smaller.  Consider the
class of all threshold functions on $\R$:
$\mathcal{H} = \{x \mapsto \ind[x \geq \theta] : \theta \in \R\}$.
This class has $\VC(\mathcal{H}) = 1$ and
$\growth_\mathcal{H}(m) = m + 1$, so the VC bound gives
$\mathcal{R}_m \leq \sqrt{2\log(em)/m}$.  But on a sample where all
points lie in $\{0,1\}$, the restriction has at most $3$ elements, and the
empirical Rademacher complexity is $O(1/\sqrt{m})$ without the
logarithmic factor.  The gain is modest here but becomes significant for
classes with distribution-dependent structure.
\end{remark}

\begin{computational}
\textbf{Rademacher complexity of linear classifiers.}
Let $\mathcal{F} = \{x \mapsto \langle w, x \rangle : \|w\|_2 \leq W\}$
and suppose $\|x_i\|_2 \leq B$ for all $i$.  Then
\[
  \Rad_S(\mathcal{F})
  = \E_\sigma\!\left[\sup_{\|w\| \leq W}
    \frac{1}{m}\left\langle w, \sum_i \sigma_i x_i\right\rangle\right]
  = \frac{W}{m}\,\E_\sigma\!\left[\left\|\sum_i \sigma_i x_i\right\|_2
    \right]
  \leq \frac{WB}{\sqrt{m}},
\]
where the last step uses
$\E[\|\sum_i \sigma_i x_i\|_2^2] = \sum_i \|x_i\|_2^2 \leq mB^2$ and
Jensen's inequality.  Note: no dependence on ambient dimension.  A
linear classifier in $\R^{10^6}$ with bounded norm has the same
Rademacher complexity as one in $\R^{10}$; only the norm and the data
radius matter.
\end{computational}


% ================================================================
% SECTION 4: THE GROWTH FUNCTION AND SAUER--SHELAH
% Brief recall (proved in Ch 10).
% ================================================================

\section{The Growth Function}
\label{sec:growth-function}

The growth function $\growth_\mathcal{H}(m) = \max_{|S|=m}
|\mathcal{H}|_S|$ mediates between VC dimension and generalization bounds.
It is the growth function, not the VC dimension directly, that enters
the uniform convergence argument (Stage~2 of the proof of
\Cref{thm:vc-uniform}).  The VC dimension is a single number; the growth
function is a whole sequence.

The Sauer--Shelah lemma (\Cref{ch:dimensions}) provides the critical
link: if $\VC(\mathcal{H}) = d$, then
$\growth_\mathcal{H}(m) \leq \sum_{i=0}^d \binom{m}{i}
\leq (em/d)^d$ for $m \geq d$.\formal{FLT_Proofs.Complexity.Rademacher.sauer_shelah_exp_bound}
The dichotomy is sharp: either $\growth_\mathcal{H}(m) = 2^m$ for all~$m$
(when $\VC = \infty$) or $\growth_\mathcal{H}(m) = O(m^d)$
(when $\VC = d < \infty$).  There is no intermediate growth rate.


% ================================================================
% SECTION 5: PAC-BAYES BOUNDS
% Profile E (obstruction).  The prior is NOT a Bayesian prior.
% ================================================================

\section{PAC-Bayes Bounds}
\label{sec:pac-bayes}

The PAC-Bayes framework replaces the worst-case uniform convergence of
Rademacher bounds with a \emph{distribution over hypotheses}.  Instead of
bounding the risk of every $h \in \mathcal{H}$ simultaneously, it bounds
the expected risk of a hypothesis drawn from a learned posterior~$\rho$,
measured relative to a data-independent prior~$\pi$.

\begin{defn}[PAC-Bayes Setting]
\label{def:pac-bayes-setting}
Fix a hypothesis space~$\mathcal{H}$, a loss
$\ell \colon \mathcal{H} \times (X \times Y) \to [0,1]$,
a \emph{prior} distribution~$\pi$ over~$\mathcal{H}$ chosen
\emph{before} seeing data, and a training sample $S \sim D^m$.
A \emph{posterior}~$\rho$ is any distribution over~$\mathcal{H}$ that
may depend on~$S$.  Define the Gibbs (mixture)
risks\formal{FLT_Proofs.Theorem.PACBayes.gibbsError}
\[
  \risk_D(\rho) = \E_{h \sim \rho}[\risk_D(h)], \qquad
  \emprisk_S(\rho) = \E_{h \sim \rho}[\emprisk_S(h)].
\]
\end{defn}

\begin{thm}[McAllester's PAC-Bayes Bound
{\cite{McAllester1999}}]
\label{thm:pac-bayes}
For any prior~$\pi$ over~$\mathcal{H}$ (chosen independently of~$S$),
any $\delta > 0$, with probability at least $1 - \delta$ over
$S \sim D^m$: for all posteriors~$\rho$ over~$\mathcal{H}$
simultaneously,
\[
  \risk_D(\rho) \leq \emprisk_S(\rho)
    + \sqrt{\frac{\KL(\rho \| \pi) + \log(2\sqrt{m}/\delta)}{2m}}.
\]
For finite~$\mathcal{H}$, the bound (with a cross-entropy complexity term in
place of $\KL(\rho\|\pi)$) is machine-checked.\formal{FLT_Proofs.Theorem.PACBayes.pac_bayes_finite}
\end{thm}

\begin{proof}
The proof has four steps: change-of-measure, moment generating function
bound, prior expectation, and optimization.

\textbf{Step 1: Change-of-measure inequality.}
For any measurable function $\phi \colon \mathcal{H} \to \R$ and
distributions $\rho$,~$\pi$, the Donsker--Varadhan variational formula
gives
\[
  \E_{h \sim \rho}[\phi(h)] \leq \KL(\rho \| \pi)
    + \log \E_{h \sim \pi}[e^{\phi(h)}].
\]
This holds for every~$\rho$; it converts a bound on the moment generating
function under~$\pi$ into a bound under any~$\rho$, at the cost of
$\KL(\rho \| \pi)$.

\textbf{Step 2: Moment generating function bound.}
Set $\phi(h) = \lambda(\risk_D(h) - \emprisk_S(h))$ for a parameter
$\lambda > 0$ to be optimized.  For a fixed~$h$ (independent of~$S$),
$\emprisk_S(h) = \frac{1}{m}\sum_{i=1}^m \ell(h, z_i)$ is an average of
i.i.d.\ $[0,1]$-valued random variables with mean~$\risk_D(h)$.  By
Hoeffding's lemma,
\[
  \E_S\!\left[e^{\lambda(\risk_D(h) - \emprisk_S(h))}\right]
  \leq e^{\lambda^2/(8m)}.
\]

\textbf{Step 3: Prior expectation.}
Since $\pi$ is independent of~$S$, we can exchange the expectations:
\[
  \E_S\!\left[\E_{h \sim \pi}\!\left[
    e^{\lambda(\risk_D(h) - \emprisk_S(h))}\right]\right]
  = \E_{h \sim \pi}\!\left[\E_S\!\left[
    e^{\lambda(\risk_D(h) - \emprisk_S(h))}\right]\right]
  \leq e^{\lambda^2/(8m)}.
\]
By Markov's inequality applied to the non-negative random variable
$\E_{h \sim \pi}[e^{\lambda(\risk_D(h) - \emprisk_S(h))}]$: with
probability $\geq 1 - \delta$,
\[
  \log \E_{h \sim \pi}\!\left[
    e^{\lambda(\risk_D(h) - \emprisk_S(h))}\right]
  \leq \frac{\lambda^2}{8m} + \log\frac{1}{\delta}.
\]

\textbf{Step 4: Combine and optimize.}
Applying Step~1 to the event from Step~3: with probability
$\geq 1 - \delta$, for all~$\rho$,
\[
  \lambda\bigl(\risk_D(\rho) - \emprisk_S(\rho)\bigr)
  \leq \KL(\rho \| \pi) + \frac{\lambda^2}{8m} + \log\frac{1}{\delta}.
\]
Dividing by~$\lambda$ and optimizing
$\lambda = \sqrt{8m(\KL(\rho\|\pi) + \log(1/\delta))}$ yields
\[
  \risk_D(\rho) - \emprisk_S(\rho)
  \leq \sqrt{\frac{2(\KL(\rho \| \pi) + \log(1/\delta))}{m}}.
\]
The displayed change-of-measure argument yields the constant
$\sqrt{2(\KL(\rho\|\pi) + \log(1/\delta))/m}$.  The sharper form stated in
\Cref{thm:pac-bayes} (with $1/(2m)$ rather than $2/m$ under the root) requires
either Catoni's thermodynamic refinement~\cite{Catoni2007} or, for
finite~$\mathcal{H}$, the direct union bound over hypotheses that the
machine-checked companion follows; we do not reproduce either sharpening here.
\end{proof}


\begin{obstbox}
\textbf{The ``Bayesian'' in PAC-Bayes is an obstruction, not a
feature.}

The PAC-Bayes bound looks Bayesian: it involves a ``prior''~$\pi$
and a ``posterior''~$\rho$, and the complexity term $\KL(\rho\|\pi)$
penalizes departure from the prior.  But the prior plays a
fundamentally different role from a Bayesian prior.

In Bayesian inference, the prior $\pi$ represents \emph{beliefs}
about which hypothesis is true.  The posterior is obtained by Bayes'
rule: $\rho(h) \propto \pi(h) \cdot \prod_i p(z_i \mid h)$.  The
prior must be honest: it should reflect actual uncertainty.

In PAC-Bayes, the prior is a \emph{reference measure} chosen for
technical convenience.  It need not represent beliefs.  The ``posterior''
is \emph{any} distribution over~$\mathcal{H}$, not necessarily the
Bayesian posterior.  The bound holds for all~$\rho$ simultaneously,
so one can optimize~$\rho$ to minimize the bound.  The optimal
$\rho$ is the \emph{Gibbs posterior}
$\rho_\lambda \propto e^{-\lambda \emprisk_S(h)} \pi(h)$, which
resembles a Bayesian posterior but with a free temperature
parameter~$\lambda$ that a true Bayesian would set to~$m$.

\emph{Type mismatch:} Bayesian prior lives in the space of beliefs;
PAC-Bayes prior lives in the space of reference measures.  They have
the same mathematical type (distributions over~$\mathcal{H}$) but
different \emph{epistemic types}.  This is why PAC-Bayes bounds hold
for \emph{any} prior, including bad ones; the bound simply becomes
large.  A Bayesian analysis with a bad prior can give arbitrarily
wrong posteriors; a PAC-Bayes analysis with a bad prior gives a
correct but loose bound.
\end{obstbox}


% ================================================================
% SECTION 6: ALGORITHMIC STABILITY
% ================================================================

\section{Algorithmic Stability}
\label{sec:stability}

Stability bounds abandon the hypothesis class entirely and focus on the
\emph{algorithm}: if replacing a single training example changes the
output hypothesis only slightly, the algorithm generalizes.

\begin{defn}[Uniform Stability {\cite{BousquetElisseeff2002}}]
\label{def:uniform-stability}
A (possibly randomized) algorithm~$A$ is \emph{$\beta$-uniformly stable}
with respect to loss~$\ell$ if, for all samples
$S = (z_1, \ldots, z_m)$ and all samples~$S^{(i)}$ obtained by
replacing~$z_i$ with an independent copy~$z_i'$:
\[
  \sup_{z \in X \times Y}
    \bigl|\ell(A(S), z) - \ell(A(S^{(i)}), z)\bigr|
  \leq \beta.
\]
\end{defn}

\begin{thm}[Stability Generalization Bound
{\cite{BousquetElisseeff2002}}]
\label{thm:stability-bound}
If $A$ is $\beta$-uniformly stable and $\ell$ takes values in $[0,1]$,
then for any $\delta > 0$, with probability at least $1 - \delta$:
\[
  |\risk_D(A(S)) - \emprisk_S(A(S))|
  \leq 2\beta + (4m\beta + 1)\sqrt{\frac{\log(1/\delta)}{2m}}.
\]
In particular, if $\beta = O(1/m)$, then the generalization gap is
$O(1/\sqrt{m})$.
\end{thm}

\begin{proof}
Define $\Delta(S) = \risk_D(A(S)) - \emprisk_S(A(S))$.  The proof
proceeds in two stages.

\emph{Stage 1: Bound $\E[\Delta(S)]$.}
Write $\risk_D(A(S)) = \E_{z'}[\ell(A(S), z')]$ and
$\emprisk_S(A(S)) = \frac{1}{m}\sum_i \ell(A(S), z_i)$.  By symmetry of
the i.i.d.\ draw,
$\E[\ell(A(S), z_i)] = \E[\ell(A(S^{(i)}), z_i')]$
up to an error of~$\beta$ (by uniform stability applied to~$S$
vs.~$S^{(i)}$).  Summing over~$i$ and using
$\E[\ell(A(S^{(i)}), z_i')] = \E[\risk_D(A(S^{(i)}))] =
\E[\risk_D(A(S))]$ (by identical distribution), one obtains
$|\E[\Delta(S)]| \leq \beta$.

\emph{Stage 2: Concentration.}
Replacing~$z_i$ changes $\emprisk_S(A(S))$ by at most $1/m$ (direct
computation) and changes $\risk_D(A(S))$ by at most~$\beta$ (via
stability).  So $\Delta(S)$ has bounded differences
$c_i \leq 2\beta + 1/m$.  McDiarmid's inequality gives
\[
  \Pr[|\Delta(S) - \E[\Delta(S)]| > t]
  \leq 2\exp\!\left(\frac{-2t^2}{\sum_i c_i^2}\right)
  \leq 2\exp\!\left(\frac{-2t^2}{m(2\beta+1/m)^2}\right).
\]
Combining with $|\E[\Delta(S)]| \leq \beta$ gives the stated bound.
\end{proof}

\begin{example}[Regularized ERM is stable]
\label{ex:regularized-erm-stable}
Consider regularized empirical risk minimization:
$A(S) = \arg\min_{h \in \mathcal{H}} \emprisk_S(h) +
\lambda \|h\|^2$.  For $\lambda$-strongly convex regularizers and
$L$-Lipschitz losses, $\beta \leq L^2/(2\lambda m)$, giving
$\beta = O(1/m)$ and hence $O(1/\sqrt{m})$
generalization~\cite{BousquetElisseeff2002}.
\end{example}

\begin{remark}[Stability beyond binary classification]
\label{rmk:stability-beyond}
A central result in stability theory is due to Shalev-Shwartz, Shamir,
Srebro, and Sridharan~\cite{ShalevShwartzShamir2010}: beyond binary
classification with 0--1 loss, \emph{learnability is equivalent to the
existence of a stable empirical risk minimizer}.  In this more general
setting (for instance, certain convex stochastic optimization
problems); uniform convergence provably fails to characterize learnability,
yet stability succeeds.  We state this without proof; the construction
separating learnability from uniform convergence is the technical core of
the cited paper.
\end{remark}


% ================================================================
% SECTION 7: INFORMATION-THEORETIC BOUNDS
% ================================================================

\section{Information-Theoretic Bounds}
\label{sec:info-bound}

The most recent entry in the generalization-bound landscape connects
generalization to the mutual information between the training sample and
the learned hypothesis.

\begin{thm}[Information-Theoretic Generalization Bound
{\cite{XuRaginsky2017}}]
\label{thm:xu-raginsky}
Let $A$ be a (possibly randomized) learning algorithm, $\ell$ a loss
taking values in $[0,1]$, and $W = A(S)$ the output hypothesis.  Let
$I(S; W)$ denote the mutual information between~$S$ and~$W$.  Then
\[
  \bigl|\E[\risk_D(W) - \emprisk_S(W)]\bigr|
  \leq \sqrt{\frac{I(S; W)}{2m}}.
\]
\end{thm}

\begin{proof}[Proof sketch]
The argument uses the transportation lemma, a consequence of Pinsker's
inequality and the Donsker--Varadhan representation of KL divergence.

For each index~$i$, define
$g_i(w) = \E_Z[\ell(w, Z)] - \ell(w, Z_i)$, where the first expectation
is over a fresh test point $Z \sim D$.  Then
$\E[\mathrm{gen}(W,S)] = \frac{1}{m}\sum_{i=1}^m \E[g_i(W)]$.

For each~$i$, the function $g_i$ satisfies $|g_i| \leq 1$, so $g_i(W)$
is $1$-sub-Gaussian under the marginal distribution of~$W$ given all
other samples.  The transportation lemma gives
$|\E[g_i(W)]| \leq \sqrt{2 I(W; Z_i)}$.
Averaging over~$i$ and using $\sum_i I(W; Z_i) \leq I(W; S)$ yields the
result.  This last inequality is \emph{not} the data-processing inequality:
it follows from the chain rule
$I(W; S) = \sum_i I(W; Z_i \mid Z_{<i})$ together with
$I(W; Z_i \mid Z_{<i}) \geq I(W; Z_i)$, the latter holding because the $Z_i$
are independent (equivalently, $I(Z_i; Z_{<i} \mid W) \geq 0$).
\end{proof}

\begin{remark}[Conditional mutual information]
\label{rmk:cmi}
The Xu--Raginsky bound can be vacuous for deterministic algorithms,
since $I(S; W) = H(W)$ when $W = A(S)$ is a deterministic function
of~$S$.  Steinke and Zakynthinou (2020) introduced the \emph{conditional
mutual information} (CMI) framework, which replaces $I(S; W)$ with
$I(A(\tilde{S}_U); U \mid \tilde{S})$, where $\tilde{S}$ is a
``supersample'' of $2m$ points and $U \in \{0,1\}^m$ selects which half
to train on.  The CMI is always bounded by $m \log 2$ and gives
nontrivial bounds even for deterministic algorithms.  This framework
unifies PAC-Bayes, VC, and stability bounds within a single
information-theoretic language.
\end{remark}


% ================================================================
% SECTION 8: MARGIN THEORY
% Profile F (Open Frontier).
% ================================================================

\section{Margin Theory}
\label{sec:margin}

Margin-based bounds were the original motivation for support vector
machines and remain the primary tool for analyzing neural network
generalization.

\begin{defn}[Margin]
\label{def:margin}
For a classifier $f \colon X \to \R$ and a labeled example $(x, y)$
with $y \in \{-1, +1\}$, the \emph{margin} is
$\gamma(f, x, y) = y \cdot f(x)$.  The margin is positive when the
prediction is correct, and its magnitude measures the ``confidence'' of
the prediction.
\end{defn}

\begin{thm}[Margin-Based Generalization Bound]
\label{thm:margin-bound}
Let $\mathcal{F}$ be a class of functions $f \colon X \to \R$ with
$\|f\|_\infty \leq B$, and let $\gamma > 0$.  Then for any $\delta > 0$,
with probability at least $1 - \delta$, for all $f \in \mathcal{F}$:
\[
  \Pr_{(x,y) \sim D}[y \cdot f(x) \leq 0]
  \leq \frac{1}{m}\sum_{i=1}^m \ind[y_i f(x_i) \leq \gamma]
  + \frac{4}{\gamma}\,\Rad_S(\mathcal{F})
  + \sqrt{\frac{\log(2/\delta)}{2m}}.
\]
\end{thm}

\begin{proof}[Proof sketch]
Define the ramp loss $\ell_\gamma(t) = \min(1, \max(0, 1 - t/\gamma))$,
which is $1/\gamma$-Lipschitz and satisfies
$\ind[t \leq 0] \leq \ell_\gamma(t) \leq \ind[t \leq \gamma]$.  Then
\begin{align*}
  \Pr_D[y f(x) \leq 0]
  &\leq \E_D[\ell_\gamma(y f(x))] \\
  &\leq \frac{1}{m}\sum_i \ell_\gamma(y_i f(x_i))
    + 2\Rad_S(\ell_\gamma \circ \mathcal{F})
    + \sqrt{\frac{\log(2/\delta)}{2m}} \\
  &\leq \frac{1}{m}\sum_i \ind[y_i f(x_i) \leq \gamma]
    + \frac{2}{\gamma}\cdot 2\Rad_S(\mathcal{F})
    + \sqrt{\frac{\log(2/\delta)}{2m}},
\end{align*}
using the Rademacher contraction lemma
($\Rad_S(\phi \circ \mathcal{F}) \leq L \cdot \Rad_S(\mathcal{F})$
for $L$-Lipschitz~$\phi$) in the last step.
\end{proof}

\begin{example}[Halfspaces with margin]
\label{ex:halfspace-margin}
Let $\mathcal{F} = \{x \mapsto \langle w, x \rangle : \|w\|_2 \leq W\}$
with data satisfying $\|x_i\| \leq B$.  From the computational
illustration in \Cref{sec:rademacher},
$\Rad_S(\mathcal{F}) \leq WB/\sqrt{m}$.  The margin bound gives
\[
  \Pr_D[y \langle w, x \rangle \leq 0]
  \leq \hat{L}_\gamma(w) + \frac{4WB}{\gamma\sqrt{m}}
    + \sqrt{\frac{\log(2/\delta)}{2m}},
\]
where $\hat{L}_\gamma(w)$ is the fraction of training points with margin
$\leq \gamma$.  Observe: the bound depends on $WB/\gamma$ (the ratio of
norm to margin), not on the ambient dimension.  This is why SVMs with
large-margin solutions generalize well even in very high dimensions.
\end{example}

\begin{example}[Neural networks: the open frontier]
\label{ex:neural-margin}
For a depth-$L$ network $f_W(x) = W_L \sigma(W_{L-1} \sigma(\cdots
W_1 x \cdots))$ with ReLU activation~$\sigma$, Bartlett, Foster, and
Telgarsky (2017) showed that the Rademacher complexity satisfies
\[
  \Rad_S(\mathcal{F})
  \leq \frac{B_x \cdot (\sqrt{2\log(2d_{\max})})^L
    \cdot \prod_{i=1}^L \|W_i\|_\sigma}{\sqrt{m}},
\]
where $\|W_i\|_\sigma$ is the spectral norm of the $i$-th layer and
$d_{\max}$ is the maximum width.  Combined with the margin bound, this
gives a generalization bound scaling as the product of spectral norms
divided by the margin.

The problem: for practical networks, this product is large.  A ResNet-18
on CIFAR-10 has $\prod_i \|W_i\|_\sigma \approx 10^3$, making the bound
vacuous (larger than~1).  Yet the network generalizes perfectly well.  The
resolution of this tension (finding bounds that are both non-vacuous and
predictive) remains one of the central open problems in learning theory.
\end{example}

\begin{openprob}[Tight deep network bounds]
\label{op:deep-margins}
For deep neural networks, all known generalization bounds (margin-based,
PAC-Bayes, stability, and information-theoretic) are \emph{loose} in
the following sense: the bounds scale with quantities that grow with
network size, yet empirically, larger networks generalize \emph{better}.

The specific open problem: find a generalization bound for deep networks
trained by stochastic gradient descent that (a)~is non-vacuous (smaller
than~1) on standard benchmarks, (b)~correctly predicts that increasing
width improves generalization, and (c)~does not require post-hoc
compression or quantization.  As of 2025, no known framework achieves all
three simultaneously.
\end{openprob}


% ================================================================
% SECTION 9: META-LEARNING BOUNDS
% Profile A-plus.  Brief.
% ================================================================

\section{Meta-Learning Bounds}
\label{sec:meta-pac}

\begin{thm}[Meta-PAC Bound {\cite{Baxter2000}}]
\label{thm:meta-pac}
In the meta-learning setting, the learner observes $T$ tasks, each
providing $m$ samples from a task-specific distribution~$D_t$.  The tasks
are drawn from a task distribution~$\mathcal{T}$.  Baxter's meta-PAC
bound states that if the bias class~$\mathcal{B}$ has finite complexity,
then with $T$ tasks and $m$ samples per task:
\[
  \risk_\mathcal{T}(\hat{b}) \leq \emprisk_{T,m}(\hat{b})
  + O\!\left(\sqrt{\frac{\log \mathcal{N}(\varepsilon,\mathcal{B},d)}{T}}
    + \sqrt{\frac{d_{\mathrm{task}}}{m}}\right),
\]
where $d_{\mathrm{task}}$ is the VC dimension of the task-level
hypothesis class induced by the bias.  The first term controls
meta-generalization; the second controls within-task generalization.  We
state the bound without proof; the argument combines a covering-number bound
at the meta level with a VC bound within each task~\cite{Baxter2000}.
\end{thm}


% ================================================================
% SECTION 10: THE COMPARISON
% CENTERPIECE: Profile G.  Five frameworks in parallax.
% ================================================================

\section{Five Frameworks Compared}
\label{sec:comparison}

The five frameworks bound the same quantity but use different information
and succeed in different regimes.  \Cref{tab:five-frameworks} compares
them along five axes.

\begin{table}[ht]
\centering\small
\caption{Five generalization-bound frameworks compared.  Each row
describes what the framework measures, what it bounds, where it is
tightest, and where it fails.}
\label{tab:five-frameworks}
\renewcommand{\arraystretch}{1.4}
\begin{tabularx}{\textwidth}{>{\bfseries\small}l X X X X}
\toprule
\textbf{Framework} & \textbf{Key quantity} & \textbf{What it bounds}
  & \textbf{Tight for} & \textbf{Fails on} \\
\midrule
Rademacher
  & $\Rad_S(\mathcal{F})$: noise correlation
  & Uniform over all $f \in \mathcal{F}$
  & Finite VC classes; linear models
  & Overparameterized models (vacuous) \\
\addlinespace
PAC-Bayes
  & $\KL(\rho \| \pi)$: posterior--prior divergence
  & Expected risk under posterior~$\rho$
  & Bayesian and ensemble methods
  & Requires meaningful prior \\
\addlinespace
Stability
  & $\beta$: one-sample sensitivity
  & $|\risk - \emprisk|$ for algorithm~$A$
  & Regularized ERM; SGD
  & Non-regularized ERM (unstable) \\
\addlinespace
Information
  & $I(S; W)$: mutual information
  & $|\E[\risk - \emprisk]|$
  & Noisy algorithms (DP-SGD)
  & Deterministic algorithms (vacuous) \\
\addlinespace
Margin
  & $\gamma$: prediction confidence
  & $\Pr[\text{error}]$ via margin loss
  & SVMs; kernel methods
  & Deep networks (norms too large) \\
\bottomrule
\end{tabularx}
\end{table}


\subsection{The Parallax Diagram}

\Cref{fig:five-frameworks} displays the implication and obstruction
relationships among the five frameworks.  The diagram should be read as a
\emph{parallax}: five views of the same phenomenon (generalization),
each projecting the three-way interaction between data, class, and
algorithm onto a different axis.

\begin{figure}[ht]
\centering
\begin{tikzpicture}[
    vfbox/.style={draw=defblue, line width=1pt, rounded corners=4pt,
      minimum width=2.5cm, minimum height=0.9cm, font=\small\bfseries,
      align=center, fill=defblue!10},
    litbox/.style={draw=black!45, thin, dashed, rounded corners=4pt,
      minimum width=2.5cm, minimum height=0.9cm, font=\small, align=center,
      fill=black!3},
    vedge/.style={-{Stealth[length=2.5mm]}, line width=1pt, defblue},
    ledge/.style={-{Stealth[length=2.5mm]}, semithick, black!55, densely dashed},
    cedge/.style={{Stealth[length=2.2mm]}-{Stealth[length=2.2mm]}, thick,
      black!45, dotted},
    elbl/.style={font=\scriptsize, midway, fill=white, inner sep=1.5pt},
]
% --- frameworks: verified-bound (thick blue border) vs literature-only (thin dashed) ---
\node[vfbox] (vc) at (-4.2,2.3) {VC / Uniform\\Convergence};
\node[vfbox] (rad) at (0,2.3) {Rademacher};
\node[litbox] (margin) at (4.2,2.3) {Margin};
\node[litbox] (info) at (0,0) {Information-\\Theoretic};
\node[litbox] (stab) at (-2.7,-2.5) {Stability};
\node[vfbox] (pb) at (2.7,-2.5) {PAC-Bayes};

% --- the one kernel-verified bridge: Rademacher controlled by VC via Sauer-Shelah ---
\draw[vedge] (vc) -- node[elbl] {Sauer--Shelah} (rad);

% --- literature implications / analogies ---
\draw[ledge] (margin) -- node[elbl] {specializes} (rad);
\draw[ledge] (stab) -- node[elbl] {implies} (info);
\draw[ledge, {Stealth[length=2.5mm]}-{Stealth[length=2.5mm]}]
  (info) -- node[elbl] {dual} (pb);

% --- conceptual orthogonalities (argued, not proved) ---
\draw[cedge] (rad) -- node[elbl] {orthogonal} (info);
\draw[cedge] (stab) -- node[elbl] {type mismatch} (pb);

% --- evidence-grade legend ---
\begin{scope}[shift={(-4.4,-4.0)}]
  \draw[vedge] (0,0) -- (0.7,0);
  \node[font=\scriptsize, anchor=west] at (0.82,0) {kernel-verified bound / bridge};
  \draw[ledge] (5.0,0) -- (5.7,0);
  \node[font=\scriptsize, anchor=west] at (5.82,0) {literature (cited)};
  \draw[cedge] (0,-0.52) -- (0.7,-0.52);
  \node[font=\scriptsize, anchor=west] at (0.82,-0.52) {conceptual orthogonality (argued)};
  \node[vfbox, minimum width=0.55cm, minimum height=0.3cm] at (5.27,-0.52) {};
  \node[font=\scriptsize, anchor=west] at (5.62,-0.52) {verified-bound};
  \node[litbox, minimum width=0.55cm, minimum height=0.3cm] at (8.5,-0.52) {};
  \node[font=\scriptsize, anchor=west] at (8.85,-0.52) {literature-only};
\end{scope}
\end{tikzpicture}
\caption[Six generalization frameworks graded by evidence: one risk gap, many certificates.]{The
six frameworks all bound the \emph{same} quantity, the gap between true and empirical risk, but
each reads a different slice of the data--class--algorithm interaction.  The kernel machine-checks
the bounds for the class-based frameworks: VC uniform convergence (\leanname{finite_massart_lemma}),
Rademacher complexity (\leanname{EmpiricalRademacherComplexity}), and PAC-Bayes
(\leanname{pac_bayes_finite}); these three nodes are drawn solid.  The one machine-checked
cross-framework bridge is that Rademacher complexity is controlled by VC dimension through
Sauer--Shelah (\leanname{vcdim_bounds_rademacher_quantitative}), the single solid edge.  The
remaining links are literature, proved on paper but not verified (margin specializes Rademacher,
stability implies an information bound, PAC-Bayes and the information bound are dual views), or
conceptual, argued rather than proved (information and Rademacher are orthogonal because one reads
the algorithm and the other the class; stability and PAC-Bayes are a type mismatch between
algorithm space and distribution space).  Solid marks a kernel-verified bound or bridge, a lighter
dashed style a cited literature result, and a dotted style a conceptual orthogonality that is
argued rather than proved.}
\label{fig:five-frameworks}
\end{figure}


\subsection{Cross-Framework Relationships}

The five frameworks are not independent.  Several implication, analogy,
and obstruction relationships connect them:

\begin{enumerate}[leftmargin=*, itemsep=3pt]
\item \textbf{Rademacher $\Rightarrow$ VC.}  The Rademacher complexity of
  a binary class satisfies $\mathcal{R}_m(\mathcal{H}) \leq
  \sqrt{2\VC(\mathcal{H}) \log(em/\VC(\mathcal{H}))/m}$
  (\Cref{prop:rad-vc}), so Rademacher bounds subsume and refine VC
  bounds.  The refinement is strict: Rademacher complexity adapts to the
  sample, whereas VC dimension does not.

\item \textbf{Margin $\Rightarrow$ Rademacher.}  The margin bound
  (\Cref{thm:margin-bound}) is derived \emph{from} Rademacher complexity
  applied to the ramp loss, so margin bounds are a specialization.

\item \textbf{Stability $\Rightarrow$ Information.}  A $\beta$-uniformly
  stable randomized algorithm satisfies $I(S; W) = O(m \beta^2)$.  The
  operative step is a stability-to-divergence bound (replacing one example
  perturbs the output distribution by $O(\beta)$ in KL divergence), not the
  data-processing inequality; summing the per-coordinate contributions gives
  the $m\beta^2$ scaling.  Thus, for noisy algorithms, stability bounds imply
  information-theoretic bounds.

\item \textbf{PAC-Bayes $\leftrightarrow$ Information.}  The PAC-Bayes
  bound with a data-independent prior can be interpreted as an
  information-theoretic bound: the KL divergence
  $\KL(\rho \| \pi)$ upper-bounds a specific mutual information
  term.  Conversely, the information-theoretic framework recovers
  PAC-Bayes via a change-of-measure argument.  The two are essentially
  dual views~\cite{HellstromDurisi2020}.

\item \textbf{Obstruction: Information $\not\Rightarrow$ Rademacher.}
  Information-theoretic bounds depend on the \emph{algorithm}; Rademacher
  bounds depend on the \emph{class}.  A deterministic ERM algorithm can have
  $I(S; W)$ as large as $H(W)$ (the full entropy of its output) while
  operating on a class with small Rademacher complexity.  The two frameworks
  capture orthogonal aspects of generalization.

\item \textbf{Obstruction: Stability $\not\Rightarrow$ PAC-Bayes.}
  Stability measures perturbation sensitivity of the \emph{algorithm}.
  PAC-Bayes measures divergence of the \emph{posterior} from the
  \emph{prior}.  An algorithm defines a posterior (connecting them), but
  the primary objects are different types: stability lives in the space of
  algorithms, PAC-Bayes in the space of distributions.
\end{enumerate}


\subsection{What Each Framework Sees}

\begin{figure}[ht]
\centering
\begin{tikzpicture}[
    vfbox/.style={draw=defblue, line width=1pt, rounded corners=3pt,
      minimum width=1.7cm, minimum height=0.58cm, font=\small\bfseries,
      align=center, fill=defblue!10, inner sep=2pt},
    litbox/.style={draw=black!45, thin, dashed, rounded corners=3pt,
      minimum width=1.7cm, minimum height=0.58cm, font=\small, align=center,
      fill=black!3, inner sep=2pt},
    cdot/.style={circle, fill, inner sep=2.4pt},
    clbl/.style={font=\footnotesize\bfseries},
    note/.style={font=\scriptsize, text=black!55},
]
% --- the three information sources (corners of the triangle) ---
\coordinate (D) at (0,0);
\coordinate (H) at (7,0);
\coordinate (A) at (3.5,5);
\draw[thick, fill=layergray] (D) -- (H) -- (A) -- cycle;
\node[cdot, defblue] at (D) {};
\node[cdot, thmred] at (H) {};
\node[cdot, sepgreen] at (A) {};
\node[clbl, defblue, below left=0pt and -3pt] at (D) {Data $D$};
\node[clbl, thmred, below right=0pt and -3pt] at (H) {Class $\mathcal{H}$};
\node[clbl, sepgreen, above=2pt] at (A) {Algorithm $A$};
% --- each framework placed at the source it reads; solid = verified bound, dashed = literature ---
% interior channel: Information reads the data-to-output channel I(S;W)
\node[litbox] (info) at (3.5,1.65) {Information};
\node[note, below=1pt of info] {$I(S;W)$ channel};
% class corner (right): VC and Rademacher (verified), Margin (literature, class geometry)
\node[vfbox] (vc) at (7.8,0.55) {VC};
\node[vfbox] (rad) at (7.8,1.4) {Rademacher};
\node[litbox] (margin) at (7.8,2.25) {Margin};
% algorithm-class edge (upper right): PAC-Bayes (verified, posterior over the class)
\node[vfbox] (pb) at (5.5,3.4) {PAC-Bayes};
% algorithm-data edge (upper left): Stability (literature)
\node[litbox] (stab) at (0.5,3.3) {Stability};
% --- evidence-grade legend (same convention as Fig 12.1) ---
\begin{scope}[shift={(-0.3,-1.35)}]
  \node[vfbox, minimum width=0.85cm, minimum height=0.3cm] at (0,0) {};
  \node[font=\scriptsize, anchor=west] at (0.55,0) {verified-bound};
  \node[litbox, minimum width=0.85cm, minimum height=0.3cm] at (3.6,0) {};
  \node[font=\scriptsize, anchor=west] at (4.15,0) {literature-only};
\end{scope}
\end{tikzpicture}
\caption[Generalization is read from three sources: data, class, or algorithm.]{Every
generalization bound reads one of three sources of information, the data, the hypothesis
class, or the algorithm, drawn here as the three corners of a triangle, so the choice of
framework is the choice of which source one can measure.  The class-based bounds sit at the
class corner and are machine-checked: VC uniform convergence (\leanname{finite_massart_lemma})
and Rademacher complexity (\leanname{EmpiricalRademacherComplexity}), drawn solid.  PAC-Bayes
(\leanname{pac_bayes_finite}) reads the posterior over the class induced by the algorithm and
is also machine-checked, so it sits on the algorithm--class edge and is drawn solid as well.
Stability reads the algorithm--data edge, information reads the data-to-output channel in the
interior, and margin reads the class geometry; these three are proved in the literature but
not verified, and are drawn lighter and dashed.  The obstructions of the companion
\Cref{fig:five-frameworks} are exactly these projections landing on different corners:
information and Rademacher read opposite sources, and stability and PAC-Bayes read different
edges, so the two figures are spatial complements of one another.}
\label{fig:triangle}
\end{figure}


\subsection{Which Framework When?}

No single framework dominates.  The choice depends on what is known:

\begin{itemize}[leftmargin=*, itemsep=3pt]
\item \textbf{If you know the class but not the algorithm}: use
  Rademacher or margin bounds.
\item \textbf{If you know the algorithm but not the class}: use
  stability or information-theoretic bounds.
\item \textbf{If you can choose a prior}: use PAC-Bayes.
\item \textbf{If you want algorithm-specific bounds for SGD}: stability
  and information-theoretic bounds apply; Rademacher bounds do not.
\item \textbf{If you want non-vacuous bounds for deep networks}:
  PAC-Bayes with learned priors is currently the most successful
  approach, but the bounds remain far from tight
  (\Cref{op:deep-margins}).
\end{itemize}


\begin{example}[Three bounds on the same problem]
\label{ex:three-bounds}
Consider linear classification with $\|w\| \leq 1$ on data with
$\|x\| \leq 1$ in $\R^d$, with $m$ training points and 0--1 loss.

\textbf{VC bound.}  $\VC = d+1$, giving generalization gap
$O(\sqrt{d \log m / m})$.  This scales linearly with dimension.

\textbf{Rademacher bound.}  $\Rad_S \leq 1/\sqrt{m}$
(from the computation in \Cref{sec:rademacher}), giving generalization
gap $O(1/\sqrt{m})$.  No dimension dependence.

\textbf{Margin bound.}  If the margin is $\gamma$, the bound is
$\hat{L}_\gamma + O(1/(\gamma\sqrt{m}))$.  If the data is
well-separated ($\gamma = 0.1$, $\hat{L}_\gamma = 0$), this gives
$O(10/\sqrt{m})$, a worse constant, but the zero margin-loss term
compensates.

The three bounds are complementary.  In dimension $d = 1000$ with
$m = 5000$, the VC bound gives approximately $0.6$ (nearly vacuous), the
Rademacher bound gives approximately $0.04$ (useful), and the margin
bound with $\gamma = 0.1$ gives approximately $0.14$ plus the margin
loss.
\end{example}


% ================================================================
% SECTION 11: CONCRETE COMPARISON ON NEURAL NETWORKS
% ================================================================

\section{A Case Study: Neural Networks}
\label{sec:nn-comparison}

The five frameworks yield strikingly different results on neural
networks, making them the ideal testbed for understanding what each
framework captures.

\begin{computational}
\textbf{Two-layer ReLU network on $\R^{100}$.}
Consider $f_W(x) = v^\top \sigma(Wx)$ with $W \in \R^{k \times 100}$,
$v \in \R^k$, width $k = 1000$, and ReLU activation $\sigma$.

\smallskip
\begin{tabular}{@{}lll@{}}
\toprule
\textbf{Framework} & \textbf{Bound expression} & \textbf{With $m=10^4$} \\
\midrule
VC & $O(\sqrt{k \cdot 100 \cdot \log(m) / m})$ & $\approx 3.4$ (vacuous) \\
Rademacher & $O(\|v\|_1 \|W\|_\sigma B_x / \sqrt{m})$
  & $\approx 0.3$ (if norms bounded) \\
PAC-Bayes & $O(\sqrt{\KL(\rho\|\pi)/m})$
  & $\approx 0.05$ (with good prior) \\
Stability & $O(L^2/(\lambda m))$ for reg.\ ERM
  & $\approx 0.1$ (if regularized) \\
Margin & $O(\|v\|_1 \|W\|_\sigma / (\gamma\sqrt{m}))$
  & depends on $\gamma$ \\
\bottomrule
\end{tabular}

\smallskip
The VC bound is vacuous because it scales with the parameter count
($\sim 10^5$, via the proxy $\VC = \tilde{O}(\#\text{params})$).
Rademacher and margin bounds depend on norms, which can be controlled by
regularization.  PAC-Bayes gives the tightest bound when a meaningful
prior is available (e.g., the initialization distribution of SGD).
Stability bounds require regularization or early stopping to
ensure~$\beta = O(1/m)$.  The information-theoretic bound depends on the
noise level of SGD: more noise means less $I(S;W)$ and a tighter bound,
but also worse training performance.
\end{computational}


\subsection{A Certified Bound for Attention}
\label{sec:certified-attention}

The five frameworks are theorems about hypothesis classes and algorithms in
the abstract.  For one concrete modern architecture, a softmax-attention
head, the pipeline of \Cref{sec:uniform-convergence,sec:rademacher} has been
carried out inside a proof assistant, producing a generalization guarantee
that is machine-checked down to the floating-point execution of the head.

\begin{prop}[Certified generalization for an attention head]
\label{prop:certified-attention}
Fix a context length~$n$, embedding dimension~$d$, and a bounded,
$L_\ell$-Lipschitz loss~$\ell$ with $|\ell| \leq b$.  Consider a single
attention head $f_\theta(x) = \mathrm{attn}\bigl(W_{\mathrm{dec}}(\theta);
\mathrm{prep}(x)\bigr)$ with a continuous, $L_W$-Lipschitz parameter map
$W_{\mathrm{dec}}$, over parameters $\theta$ in a Euclidean ball of
radius~$R$ in $\R^p$.  Set $\kappa = 4R\,L_\ell\,(dB)\,L_W$.  Then for every
$\varepsilon \geq 0$ and sample size~$m$, with probability at least
$1 - \exp\!\bigl(-\varepsilon^2 m / (2b^2)\bigr)$ over $S \sim P^m$,
\[
  \E_{x \sim P}[\ell(f_\theta(x))]
  \;\leq\; \frac{1}{m}\sum_{i=1}^m \ell(f_\theta(S_i))
    \;+\; \frac{24\sqrt{2}}{\sqrt{m}}
      \int_0^{2b}\!\Bigl(\sqrt{\log 2}
        + \sqrt{p\,\kappa}\;\eta^{-1/2}\Bigr)\mathrm{d}\eta
    \;+\; 2\varepsilon
    \;+\; 2 L_\ell\,\mathrm{env}(f_\theta),
\]
where $\mathrm{env}(f_\theta)$ is the floating-point execution envelope of
the head.
\end{prop}

\begin{proof}[Proof idea]
The bound instantiates the chapter's machinery on the attention head viewed
as a real-valued function class indexed by~$\theta$.  The middle term is the
\emph{Dudley entropy integral} (\Cref{ch:compression}) of the
parameter-perturbation cover of the head: the Lipschitz constant $\kappa$
controls the covering numbers, and the integral $\int_0^{2b}\sqrt{\log
\mathcal{N}(\eta)}\,\mathrm{d}\eta$ takes the displayed closed form because a
radius-$R$ ball in $\R^p$ has metric entropy $O(p\log(1/\eta))$.  The
high-probability lift is the bounded-differences (McDiarmid) step of
\Cref{thm:rademacher-bound}, with per-example increment $2b/m$.  The one new
ingredient is the rounding term: the verified statement bounds the
\emph{executed}, finite-precision head, so it carries an explicit envelope
$\mathrm{env}(f_\theta)$ separating the idealized map from its floating-point
realization.  The full derivation is
formalized.\formal{TLT.attnHead_certified_generalization}
\end{proof}

The construction lifts from one head to a depth-graded multi-head encoder
stack: the budget becomes the Dudley capacity of the depth-replicated block,
growing through the per-layer perturbation envelope with depth and head
count, and the resulting predicate (true risk exceeds empirical risk plus
budget plus rounding only on a $\delta$-event) is discharged for the stack
indexed by its architectural
configuration.\formal{TLT.config_capacity_law}  That predicate is itself the
machine-checked definition of a generalization
certificate.\formal{TLT.CertifiedGeneralization}

\begin{remark}[What the certificate settles, and what it does not]
\label{rmk:certified-attention}
\Cref{prop:certified-attention} is an end-to-end verified generalization
bound for a transformer component: the abstract theory of this chapter,
instantiated and checked for attention.  It does \emph{not} resolve
\Cref{op:deep-margins}.  The Dudley budget still grows with the parameter
radius~$R$, the depth, and the head count, so on a large pretrained model the
certificate is correct but loose, exactly as the classical bounds are.
What it adds is a different axis of progress: a \emph{trustworthy} bound,
whose constants and finite-precision corrections have been audited by a
kernel rather than estimated by hand.
\end{remark}


% ================================================================
% OPEN PROBLEMS
% ================================================================

\section{Open Problems}
\label{sec:generalization-open}

The deep-network gap of \Cref{op:deep-margins} is the headline question, but
the chapter's five frameworks and their machine-checked fragments open a
wider frontier.  Each problem below is tagged as a \emph{known unknown} (KU),
sharply posed but unresolved, or an \emph{unknown known} (UK), where the
right invariant or formal object is still missing; a design-lab or literature
anchor is named where one exists.

\begin{openprob}[Formalizing the data-dependent Rademacher bound]
\label{op:formalize-rademacher}
The empirical Rademacher complexity, the Massart finite lemma, and the VC
bound $\mathcal{R}_m \leq \sqrt{2d\log(em/d)/m}$ are machine-checked
(\Cref{def:rademacher,prop:rad-vc}), but the assembled high-probability bound
of \Cref{thm:rademacher-bound} (symmetrization plus the two McDiarmid
applications) is so far only a placeholder in the companion development.  Can
the full data-dependent bound be formalized end to end?  A known unknown:
classical on paper~\cite{BartlettMendelson2002}, open as a verified artifact.
\end{openprob}

\begin{openprob}[A non-vacuous certified bound for attention]
\label{op:tighten-certified-attention}
The certificate of \Cref{prop:certified-attention} is correct but its Dudley
budget grows with the parameter radius, depth, and head count
(\leanname{config_capacity_law}).  Can one machine-check a bound for a
realistic attention stack that is also \emph{non-vacuous} on a benchmark: a
verified analogue of the PAC-Bayes results for deep networks?  A known
unknown, and the verified counterpart of \Cref{op:deep-margins}.
\end{openprob}

\begin{openprob}[Discharging the expressivity leg]
\label{op:expressivity-leg}
The transformer design law assembles a certified capacity bound but states an
\emph{expressivity} obligation as an open hypothesis: that the stack can
realize a target risk at all.  Can the expressivity leg be discharged with
explicit width/depth, turning the design law into an unconditional
generalization-and-approximation guarantee?  An unknown known: the
approximation invariant for the certified stack is not yet
formalized~\cite{XuRaginsky2017}.
\end{openprob}

\begin{openprob}[Making the CMI unification precise]
\label{op:cmi-unification}
The conditional-mutual-information framework (\Cref{rmk:cmi}) is said to
unify PAC-Bayes, VC, and stability bounds.  Is there a single CMI inequality
from which all three are recovered as corollaries, with explicit constants?
An unknown known: the unification is asserted but no master theorem with the
three recoveries written out exists~\cite{SteinkeZakynthinou2020}.
\end{openprob}

\begin{openprob}[When is uniform convergence necessary?]
\label{op:uniform-necessity}
Shalev-Shwartz et al.\ show that beyond binary classification, learnability
can hold without uniform convergence (\Cref{rmk:stability-beyond}).  Exactly
which loss/hypothesis structures force uniform convergence to be necessary,
not merely sufficient?  A known unknown: the boundary between the two regimes
is uncharted~\cite{ShalevShwartzShamir2010}.
\end{openprob}

\begin{openprob}[Generalization of interpolators]
\label{op:benign-overfitting}
Modern networks interpolate the training data ($\emprisk_S = 0$) yet
generalize, contradicting the intuition behind every bound in this chapter.
Is there a complexity measure that is small precisely for the interpolators
that generalize (benign overfitting) and large for those that do not (the
double-descent peak)?  A known unknown driving current theory.
\end{openprob}

\begin{openprob}[Optimal data-dependent priors]
\label{op:pac-bayes-prior}
PAC-Bayes is tightest with a prior close to the learned posterior, but the
prior must be data-independent.  Is there a principled, optimal way to build
a data-dependent prior (via data splitting or differential privacy) that
provably minimizes the bound?  A known unknown opened by the non-vacuous
deep-network results~\cite{DziugaiteRoy2017}.
\end{openprob}

\begin{openprob}[A formal margin bound for attention]
\label{op:margin-attention}
The spectral-norm margin bound for deep networks (\Cref{ex:neural-margin})
has no verified counterpart for attention.  Can one machine-check a
margin-based bound for a softmax-attention class, with the product of
query/key/value operator norms in place of the layer spectral
norms~\cite{BartlettFosterTelgarsky2017}?  A known unknown on the transformer
side, complementary to \Cref{prop:certified-attention}.
\end{openprob}

\begin{openprob}[Rademacher lower bounds]
\label{op:rademacher-lower}
Every bound here is an upper bound on the generalization gap.  For which
classes is the Rademacher complexity also a \emph{lower} bound on the
worst-case gap, so that $\mathcal{R}_m$ exactly captures learnability rather
than merely bounding it?  An unknown known: a matching lower-bound invariant
is missing for general classes.
\end{openprob}

\begin{openprob}[Non-vacuous information bounds for deterministic learners]
\label{op:info-deterministic}
The Xu--Raginsky bound is vacuous when $W = A(S)$ is deterministic, since
$I(S;W) = H(W)$ (\Cref{rmk:cmi}).  Apart from the CMI construction, is there a
generalization bound for deterministic algorithms phrased in genuinely
information-theoretic quantities that stays finite?  A known unknown at the
frontier of the information-theoretic tradition.
\end{openprob}


% ================================================================
% SECTION 12: HISTORICAL CONTEXT
% ================================================================

\begin{historical}
The study of generalization bounds spans five decades and five research
communities.

The VC theory (Vapnik--Chervonenkis, 1971) established the first uniform
convergence bounds.  The connection to PAC learning was made explicit by
Blumer et al.\ (1989).  Rademacher complexity was developed independently
by Koltchinskii (2001) and Bartlett--Mendelson (2002) as a data-dependent
refinement.

PAC-Bayes bounds originated with Shawe-Taylor and Williamson (1997) and
were formalized by McAllester (1998, 1999).  Catoni (2007) developed the
optimal ``thermodynamic'' form.  The framework gained renewed interest
around 2017 when Dziugaite and Roy showed that PAC-Bayes bounds can be
made non-vacuous for deep networks.

Algorithmic stability was introduced by Devroye and Wagner (1979) in a
simpler form, and the modern uniform stability framework was established
by Bousquet and Elisseeff (2002).  The connection to learnability beyond
binary classification was established by Shalev-Shwartz et al.\ (2010).

Information-theoretic bounds are the most recent tradition, beginning
with Russo and Zou (2016) and Xu and Raginsky (2017).  The conditional
mutual information framework of Steinke and Zakynthinou (2020) provided
the unifying perspective.

Margin theory has the longest applied history (Vapnik, 1995; Bartlett,
1998) and was reinvigorated by Bartlett, Foster, and Telgarsky (2017)
for deep networks.

The unification of these traditions (showing they are all projections
of a single underlying phenomenon) remains an active research
direction.
\end{historical}


% ================================================================
% NOTES
% ================================================================

\bigskip

\begin{remark}[The chapter in context]
This chapter completes Part~III's treatment of quantitative measures of
learning complexity.  \Cref{ch:dimensions} developed the combinatorial
dimensions (VC, Littlestone, DS).  \Cref{ch:compression} connected
learnability to sample compression.  The present chapter adds the
\emph{analytic} layer: bounds that read the algorithm, the data, and the
geometry of the learned hypothesis on top of the class itself.
The five frameworks exemplify a recurring theme of this book: the same
mathematical phenomenon (in this case, the gap between training and
test performance) looks different depending on which structural
information one chooses to measure.  No single framework captures the
full picture, and the obstructions between frameworks (\Cref{fig:five-frameworks})
are as informative as the implications.
\end{remark}

\subsection*{Counterfactual exercises: generalizing Figure~\ref{fig:five-frameworks}}

The parallax of \Cref{fig:five-frameworks} is unusually lopsided as figures in this book go: most of its content is a frontier.  Of six framework nodes only three are drawn solid (VC uniform convergence \leanname{finite_massart_lemma}, Rademacher complexity \leanname{EmpiricalRademacherComplexity}, and PAC-Bayes \leanname{pac_bayes_finite}, all class-based); the other three (stability, information, margin) are literature.  Of six cross-framework links exactly one is solid, the Sauer--Shelah bridge by which VC dimension controls Rademacher complexity \leanname{vcdim_bounds_rademacher_quantitative}; the rest are either literature implications (margin specializing Rademacher, stability implying an information bound, PAC-Bayes and information being dual) or merely-argued conceptual edges (the algorithm-versus-class orthogonality and the stability-versus-PAC-Bayes type mismatch).  So when each exercise below perturbs the picture, ask which of three things the perturbation produces: an \emph{instance} that the one verified core or a single graded edge already exhibits, a \emph{boundary} where toggling one node sort, edge grade, or information axis flips the evidence grade or the direction, or a \emph{frontier} where a literature or conceptual edge still waits to be promoted to verified.  Most of these land on the frontier, which is the honest shape of the subject.

\begin{enumerate}[leftmargin=*,itemsep=3pt]
\item \textbf{The three solid nodes are exactly the class-based bounds the kernel checks.} The figure's three solid borders are a claim about evidence, and this is the defining verified instance of the whole diagram: VC uniform convergence, Rademacher complexity, and PAC-Bayes are the frameworks whose bounds are machine-checked (\cref{thm:vc-uniform}, \leanname{finite_massart_lemma}; \cref{def:rademacher}, \leanname{EmpiricalRademacherComplexity}; \cref{thm:pac-bayes}, \leanname{pac_bayes_finite}), so a solid border asserts exactly \emph{that the bound is a kernel theorem}. The three lighter-bordered nodes, stability, information, and margin, are drawn weaker precisely because their bounds (\cref{thm:stability-bound}, \cref{thm:xu-raginsky}, \cref{thm:margin-bound}) are literature with no verified handle.

\item \textbf{The one solid edge is the only machine-checked cross-framework bridge.} Take the single solid arrow between Rademacher and VC on its own. It is the lone settled link in the diagram: Rademacher complexity is bounded by VC dimension through Sauer--Shelah, $\mathcal{R}_m(\mathcal{H}) \le \sqrt{2\VC(\mathcal{H})\log(em/\VC(\mathcal{H}))/m}$ (\cref{prop:rad-vc}, \leanname{vcdim_bounds_rademacher_quantitative}), the one cross-framework relationship the kernel certifies. The old figure's separate ``refines'' arrow and dotted ``Sauer--Shelah'' arrow named the \emph{same} fact, which is why the re-graded figure collapses them into one solid edge and leaves nothing else solid between frameworks. This is the figure's instance.

\item \textbf{The parallax is one risk-gap quantity seen from six sides.} Read the whole figure as six projections of one object rather than six theories, and it becomes the organizing instance: every framework bounds the single quantity $|\risk - \emprisk|$ (or its expectation), and the rows of \cref{tab:five-frameworks} differ only in the \emph{key quantity} each reads (noise correlation, posterior--prior divergence, one-sample sensitivity, mutual information, margin). The diagram is the abstract witness that generalization is one quantity admitting many certificates, three of them, the class-based ones, machine-checked, and the rest borrowed from the literature.

\item \textbf{PAC-Bayes is the verified bound that lives in distribution space.} The PAC-Bayes node anchors the posterior axis: the McAllester finite bound on the Gibbs risk is machine-checked (\cref{thm:pac-bayes}, \leanname{pac_bayes_finite}, with \leanname{gibbsError} the verified mean error of the posterior-sampled predictor), so PAC-Bayes joins VC and Rademacher among the solid nodes even though it reads the \emph{posterior} rather than the class. That same kernel handle is what keeps the conceptual stability-to-PAC-Bayes edge a type relation and not a theorem, since stability reads the algorithm and PAC-Bayes the distribution it induces. This is the instance that fixes the posterior reading.

\item \textbf{The two obstructions are orthogonal information sources, not separations.} Toggle the figure's reading of the conceptual obstructions from ``proved gap'' to ``different information axis,'' and the figure's principal boundary appears: information-theoretic bounds depend on the \emph{algorithm} and Rademacher bounds on the \emph{class}, so a deterministic ERM can carry $I(S;W)=H(W)$ on a class of small Rademacher complexity (\cref{sec:comparison}), and stability lives in algorithm space while PAC-Bayes lives in distribution space. Neither obstruction is a separation theorem nor even a formal literature separation, so both are drawn in the conceptual style, argued not proved, and reversing either direction does not turn it into a verified edge.

\item \textbf{Margin is a Rademacher specialization, and that specialization is literature.} The margin-to-Rademacher edge is the boundary that fixes the edge's grade: the margin bound is \emph{derived from} Rademacher complexity applied to the ramp loss (\cref{thm:margin-bound}), so the arrow points from margin into Rademacher as a specialization, yet because thm:margin-bound is not machine-checked the edge is drawn lighter than the solid rad--vc bridge. Promoting it would require a verified ramp-loss contraction, exactly the kind of step still missing for the assembled bounds.

\item \textbf{The stability-to-information edge is one direction only, and unverified.} Follow the stability-to-information arrow and ask whether the reverse is drawn; it is not, and that asymmetry is the boundary between a literature implication and an absent one. A $\beta$-uniformly stable randomized algorithm satisfies $I(S;W)=O(m\beta^2)$ (\cref{sec:comparison}, \cref{thm:stability-bound}), a one-way implication with no kernel handle, so the figure draws stability into information in the literature style and leaves the converse undrawn. Both endpoints are lighter-bordered nodes, which confirms that an edge between two unverified frameworks cannot be solid even when the implication is classical.

\item \textbf{The Rademacher generalization bound is the figure's nearest open frontier.} Ask what it would take to draw a \emph{second} solid edge out of the Rademacher node, the full generalization bound rather than the VC bridge: this is the figure's nearest verification frontier. The empirical Rademacher complexity, the Massart lemma, and the VC bound are machine-checked, but the assembled high-probability bound of \cref{thm:rademacher-bound} (symmetrization plus the two McDiarmid steps) is still a placeholder (\cref{op:formalize-rademacher}, \cite{BartlettMendelson2002}), so the Rademacher node is solid for its \emph{complexity measure} while its generalization \emph{guarantee} stays literature. Closing this would promote Rademacher's own bound to solid, alongside its already-solid bridge to VC.

\item \textbf{Formalizing stability and information would flip two lighter nodes to solid.} Which perturbation promotes the algorithm-based frameworks? Here is the frontier that the conceptual obstructions hide: the stability bound (\cref{thm:stability-bound}) and the Xu--Raginsky information bound (\cref{thm:xu-raginsky}, \cite{XuRaginsky2017}) have no verified handle, and a single conditional-mutual-information master inequality is conjectured to recover VC, PAC-Bayes, and stability at once (\cref{op:cmi-unification}, \cite{SteinkeZakynthinou2020}). Proving and formalizing it would recolor the stability and information nodes solid and replace the stability-to-information literature edge by a verified one, leaving only the genuinely orthogonal algorithm-versus-class edge conceptual.

\item \textbf{Is every framework a section of one risk-gap functor over its own fibration?} Read ``one quantity, many certificates'' as a universal property, and you reach the figure's deepest frontier. Each framework would be a section of a single risk-gap functor over a different fibration of the learning problem, the class for VC and Rademacher, the algorithm for stability and information, the posterior for PAC-Bayes, the operator norms for margin, so the verified core would read ``the functor admits a section over the class fibration'' (the three solid nodes) and the conceptual obstructions would read ``the class and algorithm fibrations do not refine one another,'' explaining the orthogonality as a fact about base categories rather than a separation. Whether such a single generalization-certificate object exists, of which the certified attention bound (\cref{prop:certified-attention}, \leanname{CertifiedGeneralization}) is a sixth, algorithm-specific section, or the fibrations are provably non-unifiable, is open (\cref{op:cmi-unification}, \cref{op:tighten-certified-attention}), so this reading is drawn dashed.
\end{enumerate}

\subsection*{Counterfactual exercises: generalizing Figure~\ref{fig:triangle}}

Where \Cref{fig:five-frameworks} graded the links, \Cref{fig:triangle} grades the geometry.  Its grammar is the placement: generalization information has three sources, the data, the hypothesis class, and the algorithm, set at the corners of a triangle, and every bound is a projection that reads one corner or one edge.  Three framework nodes are solid because their bounds are machine-checked (VC and Rademacher at the class corner, \leanname{finite_massart_lemma} and \leanname{EmpiricalRademacherComplexity}; PAC-Bayes on the algorithm--class edge, \leanname{pac_bayes_finite}); three are lighter and dashed because their bounds are literature (stability on the algorithm--data edge, information in the interior channel, margin on the class geometry).  No edge here is a theorem; the figure is pure projection, and the only verified content it carries is the grade of each node.  The obstructions of \Cref{fig:five-frameworks} reappear here as projections that simply land on different corners.  So as each exercise below perturbs the projection, classify what it yields: an \emph{instance} whose corner placement the verified core or the literature already exhibits, a \emph{boundary} where toggling one node grade or information source flips the evidence grade or the source being read, or a \emph{frontier} where a literature reading still waits to be promoted to a verified source.

\begin{enumerate}[leftmargin=*,itemsep=3pt]
\item \textbf{The class corner carries both verified bounds.} Take the two solid nodes at the class corner together and you have the figure's defining verified instance: VC uniform convergence and Rademacher complexity both read the hypothesis class and nothing else, and both have machine-checked bounds (\cref{thm:vc-uniform}, \leanname{finite_massart_lemma} and \leanname{sauer_shelah_exp_bound}; \cref{def:rademacher}, \leanname{EmpiricalRademacherComplexity} and \leanname{vcdim_bounds_rademacher_quantitative}). The class corner is therefore the most fully certified source in the diagram, and the solid border at that corner asserts exactly that the class-source reading is a kernel theorem.

\item \textbf{PAC-Bayes is the verified reading on the algorithm--class edge.} Why is the PAC-Bayes node solid yet off the class corner? This is the instance that anchors the posterior reading: the McAllester finite bound on the Gibbs risk is machine-checked (\cref{thm:pac-bayes}, \leanname{pac_bayes_finite}, with \leanname{gibbsError} the verified mean error of the posterior-sampled predictor), and PAC-Bayes reads the posterior over the class \emph{induced by the algorithm}, so its node sits on the edge joining the algorithm and class corners, solid like VC and Rademacher but reading two sources at once.

\item \textbf{One problem read from three corners is the three-bounds example.} \Cref{ex:three-bounds} is a tour of the triangle, and so the figure's organizing instance: linear classification with bounded norm admits a VC bound (class corner, $\VC=d+1$), a Rademacher bound (class corner, $\Rad_S\le 1/\sqrt{m}$), and a margin bound (class geometry), all on the \emph{same} problem. The diagram is the abstract witness that a single learning problem is bounded by reading whichever corner one can measure, and the three numbers in the example differ only because they project different sources.

\item \textbf{Which framework when is a corner-selection rule.} The ``Which Framework When?'' guide is a map back onto the triangle, which makes it the instance that names the figure's use: knowing the class but not the algorithm selects the class corner (Rademacher or margin), knowing the algorithm but not the class selects the algorithm readings (stability or information), and being able to choose a prior selects the algorithm--class edge (PAC-Bayes). The choice of framework is literally the choice of which corner of \cref{fig:triangle} one can measure.

\item \textbf{The obstructions are different-corner projections, not separations.} Toggle the companion figure's two conceptual edges onto this triangle and the figure's principal boundary appears: information reads the data-to-output channel in the interior while Rademacher reads the class corner, so the \cref{fig:five-frameworks} ``orthogonal'' edge is just these two nodes sitting at different places (\cref{sec:comparison}), and stability reads the algorithm--data edge while PAC-Bayes reads the algorithm--class edge, so the ``type mismatch'' edge is two nodes on different edges sharing only the algorithm corner. Neither is a separation theorem, which is why both figures draw them without a verified arrow.

\item \textbf{Stability and information read the algorithm, not the class.} Why are the two interior or algorithm-edge nodes dashed? This is the boundary that fixes their grade: stability measures the algorithm's sensitivity to one data point (\cref{thm:stability-bound}) and information measures the mutual information $I(S;W)$ between the sample and the algorithm's output (\cref{thm:xu-raginsky}, \cite{XuRaginsky2017}), so both read the algorithm rather than the class and neither has a verified handle, hence both nodes are drawn lighter. A deterministic ERM can carry $I(S;W)=H(W)$ on a class of small Rademacher complexity, which confirms the interior reading is genuinely a different source from the class corner.

\item \textbf{Margin reads class geometry, and that reading is literature.} The margin node near the class--algorithm side marks the boundary between a class-adjacent reading and a verified one: the margin bound is derived from Rademacher complexity applied to the ramp loss (\cref{thm:margin-bound}), so it reads the class geometry together with the predictor's margin, yet because thm:margin-bound is not machine-checked the node is drawn dashed even though it sits beside the solid class corner. Promoting it would require a verified ramp-loss contraction, the same step missing for the assembled Rademacher bound.

\item \textbf{Formalizing the algorithm readings would recolor the interior.} Which perturbation turns the dashed interior and algorithm-edge nodes solid? Here is the frontier the dashed nodes mark: the stability bound (\cref{thm:stability-bound}) and the Xu--Raginsky information bound (\cref{thm:xu-raginsky}, \cite{XuRaginsky2017}) have no verified handle, and the Rademacher node is solid for its complexity measure while its assembled high-probability \emph{guarantee} (\cref{thm:rademacher-bound}) is still a placeholder (\cref{op:formalize-rademacher}, \cite{BartlettMendelson2002}). Formalizing these would recolor the algorithm-reading and interior nodes solid and certify sources beyond the class corner.

\item \textbf{Is there a single bound that reads all three corners at once?} Could one certificate project the whole triangle rather than one corner or edge? That is the figure's unifying frontier: a single conditional-mutual-information master inequality is conjectured to recover VC, PAC-Bayes, and stability simultaneously (\cref{op:cmi-unification}, \cite{SteinkeZakynthinou2020}), which would be one bound reading the class, the algorithm, and the posterior together, collapsing the three solid-and-dashed readings into a single source-spanning certificate. Until it is proved and formalized the corners remain separate projections.

\item \textbf{Is the triangle the base of a fibration with each framework a section?} Read ``three sources, many projections'' as a universal property and you reach the figure's deepest frontier. Each framework would be a section of one risk-gap functor over a different face of the learning problem, the class face for VC and Rademacher, the algorithm face for stability and information, the posterior for PAC-Bayes, the operator norms for margin, so the verified core would read ``a section exists over the class face'' (the solid corner) and the obstructions would read ``the class and algorithm faces do not refine one another,'' explaining the orthogonality as a fact about base faces rather than a separation. Whether such a single generalization-certificate object exists, of which the certified attention bound (\cref{prop:certified-attention}, \leanname{CertifiedGeneralization}) is an algorithm-face section, or the faces are provably non-unifiable, is open (\cref{op:cmi-unification}, \cref{op:tighten-certified-attention}), so this reading is drawn dashed.
\end{enumerate}
